\documentclass[11pt]{article}

\usepackage[a4paper,margin=1in]{geometry}
\usepackage[T1]{fontenc}
\usepackage[utf8]{inputenc}
\usepackage{lmodern}
\usepackage{microtype}
\usepackage{amsmath,amssymb,amsthm,mathtools}
\usepackage{bbm}
\usepackage{booktabs}
\usepackage{array}
\usepackage{enumitem}
\usepackage[hidelinks]{hyperref}

\title{Wasserstein Polyak--{\L}ojasiewicz Inequalities,\\ Linear Convergence, and Functional Inequality Examples}
\author{}
\date{\today}

\newtheorem{theorem}{Theorem}[section]
\newtheorem{proposition}[theorem]{Proposition}
\newtheorem{lemma}[theorem]{Lemma}
\newtheorem{corollary}[theorem]{Corollary}
\theoremstyle{definition}
\newtheorem{definition}[theorem]{Definition}
\newtheorem{example}[theorem]{Example}
\newtheorem{remark}[theorem]{Remark}

\DeclareMathOperator*{\argmin}{arg\,min}
\DeclareMathOperator{\Ent}{Ent}
\DeclareMathOperator{\Var}{Var}
\DeclareMathOperator{\KL}{KL}
\newcommand{\Ptwo}{\mathcal P_2}
\newcommand{\F}{\mathcal F}
\newcommand{\G}{\mathcal G}
\newcommand{\D}{\mathcal D}
\newcommand{\I}{\mathcal I}
\newcommand{\W}{W_2}
\newcommand{\R}{\mathbb R}
\newcommand{\dd}{\,\mathrm d}
\newcommand{\gradW}{\nabla_W}
\newcommand{\slope}{|\partial \F|}
\newcommand{\rhostar}{\rho_\star}
\newcommand{\distW}{\operatorname{dist}_{\W}}

\begin{document}
\maketitle

\begin{abstract}
This note records a useful dictionary: a Wasserstein Polyak--{\L}ojasiewicz inequality is exactly an entropy--entropy production inequality for the chosen Wasserstein energy. Under the usual energy dissipation identity for the Wasserstein gradient flow, it gives exponential decay of the energy gap and, by a tail-length argument, exponential convergence in Wasserstein distance to a minimizer selected by the flow. The exponent is $2\kappa$ for the energy gap and $\kappa$ for the $W_2$ distance itself. Positive Wasserstein geodesic convexity implies the Wasserstein--PL inequality, but the converse is false. A concrete counterexample is given by the potential energy generated by the nonconvex one-dimensional potential
\[
        V(x)=x^2+2\sin^2 x.
\]
The last section lists several classical and nonclassical functional inequalities that are Wasserstein--PL inequalities for particular energies, and explains how the distance estimate specializes to Talagrand-type transportation inequalities in entropy examples.
\end{abstract}

\tableofcontents

\section{Setting and notation}

Let $M$ be either $\R^d$ or a smooth complete Riemannian manifold, and let
$\Ptwo(M)$ denote the probability measures with finite second moment, endowed with the $2$-Wasserstein distance $\W$.
For a functional
\[
        \F:\Ptwo(M)\to (-\infty,+\infty]
\]
we write
\[
        \F_\star := \inf_{\rho\in\Ptwo(M)}\F(\rho),
        \qquad
        \operatorname{Argmin}\F := \{\rho:\F(\rho)=\F_\star\}.
\]
The metric slope of $\F$ at $\rho$ is
\begin{equation}
        |\partial \F|(\rho)
        :=
        \limsup_{\eta\to \rho}
        \frac{\bigl(\F(\rho)-\F(\eta)\bigr)_+}{\W(\rho,\eta)}.
        \label{eq:metric-slope}
\end{equation}
When $\rho$ has a density and $\F$ has first variation $\delta \F/\delta\rho$, the formal Otto calculus gives
\begin{equation}
        |\partial \F|^2(\rho)
        = \|\gradW \F(\rho)\|_{T_\rho\Ptwo}^2
        = \int_M \left|\nabla \frac{\delta \F}{\delta \rho}(\rho)(x)\right|^2\dd \rho(x).
        \label{eq:slope-first-variation}
\end{equation}
Constants in $\delta\F/\delta\rho$ do not matter, because only its gradient appears.

The Wasserstein gradient flow of $\F$ is formally
\begin{equation}
        \partial_t\rho_t
        = \nabla\cdot\left(\rho_t\nabla\frac{\delta \F}{\delta\rho}(\rho_t)\right).
        \label{eq:wgf-pde}
\end{equation}
For smooth solutions, and more generally for curves of maximal slope when the slope is a strong upper gradient, one has the energy dissipation identity
\begin{equation}
        \frac{\dd}{\dd t}\F(\rho_t)
        = - |\partial \F|^2(\rho_t)
        = -\|\dot\rho_t\|_{T_{\rho_t}\Ptwo}^2
        \quad\text{for a.e. }t>0.
        \label{eq:edi}
\end{equation}
The arguments below use only \eqref{eq:edi} and the functional inequalities stated explicitly.

\section{Wasserstein--PL implies energy and Wasserstein-distance convergence}

\begin{definition}[Wasserstein Polyak--{\L}ojasiewicz inequality]
Let $\kappa>0$. We say that $\F$ satisfies the $\W$-PL inequality with constant $\kappa$ if
\begin{equation}
        |\partial \F|^2(\rho)
        \ge 2\kappa\bigl(\F(\rho)-\F_\star\bigr)
        \quad\text{for every }\rho\in \operatorname{Dom}\F.
        \label{eq:wpl}
\end{equation}
Equivalently, in the smooth Otto notation,
\begin{equation}
        \int_M \left|\nabla \frac{\delta \F}{\delta \rho}(\rho)\right|^2\dd \rho
        \ge 2\kappa\bigl(\F(\rho)-\F_\star\bigr).
        \label{eq:wpl-otto}
\end{equation}
\end{definition}

The factor $2$ is a convention. With this convention, the logarithmic Sobolev inequality
$2\kappa\,\Ent_\pi(\rho)\le I_\pi(\rho)$ is a $\W$-PL inequality with constant $\kappa$.


\begin{theorem}[Energy, finite length, and Wasserstein-distance convergence]\label{thm:wpl-flow-convergence}
Assume that $\F_\star> -\infty$ and that $(\rho_t)_{t\ge0}$ is a global Wasserstein gradient flow satisfying the full energy dissipation identity
\begin{equation}
        \frac{{\rm d}}{{\rm d}t}\F(\rho_t)
        =-|\dot\rho_t|^2
        =-|\partial\F|^2(\rho_t)
        \quad\text{for a.e. }t>0.
        \label{eq:full-ede}
\end{equation}
Assume also that $\F$ satisfies the $\W$-PL inequality \eqref{eq:wpl}. Put
\[
        E(t):=\F(\rho_t)-\F_\star.
\]
Then, for $0\le s\le t$,
\begin{equation}
        E(t)
        \le e^{-2\kappa(t-s)}E(s),
        \label{eq:energy-rate-from-s}
\end{equation}
and the length of the flow segment satisfies
\begin{equation}
        \operatorname{Length}_{\W}\bigl((\rho_r)_{r\in[s,t]}\bigr)
        \le
        \sqrt{\frac{2}{\kappa}}\bigl(\sqrt{E(s)}-\sqrt{E(t)}\bigr).
        \label{eq:length-bound-finite-time}
\end{equation}
In particular,
\begin{equation}
        \W(\rho_s,\rho_t)
        \le
        \sqrt{\frac{2}{\kappa}}\bigl(\sqrt{E(s)}-\sqrt{E(t)}\bigr).
        \label{eq:w-between-times}
\end{equation}
If, in addition, $\Ptwo(M)$ is complete along the considered curve and $\F$ is lower semicontinuous with respect to $\W$, then $\rho_t$ converges in $\W$ to a minimizer $\rho_\infty\in\operatorname{Argmin}\F$ and
\begin{equation}
        \W(\rho_t,\rho_\infty)
        \le
        \sqrt{\frac{2}{\kappa}}\sqrt{\F(\rho_t)-\F_\star}
        \le
        \sqrt{\frac{2}{\kappa}}\sqrt{\F(\rho_0)-\F_\star}\,e^{-\kappa t}.
        \label{eq:w-to-limit-rate}
\end{equation}
Consequently,
\begin{equation}
        \distW(\rho_t,\operatorname{Argmin}\F)
        \le
        \sqrt{\frac{2}{\kappa}}\sqrt{\F(\rho_0)-\F_\star}\,e^{-\kappa t}.
        \label{eq:w-to-argmin-rate}
\end{equation}
\end{theorem}

\begin{proof}
The energy estimate is the usual Gronwall argument: by \eqref{eq:full-ede} and \eqref{eq:wpl},
\[
        E'(t)=-|\partial\F|^2(\rho_t)\le -2\kappa E(t)
        \quad\text{for a.e. }t.
\]
This gives \eqref{eq:energy-rate-from-s}.

The distance estimate uses one extra observation: PL bounds the inverse speed along the dissipating curve. On the set where $E(r)>0$, \eqref{eq:wpl} gives
\[
        |\partial\F|(\rho_r)
        \ge \sqrt{2\kappa E(r)}.
\]
Using \eqref{eq:full-ede},
\begin{align*}
        \operatorname{Length}_{\W}\bigl((\rho_r)_{r\in[s,t]}\bigr)
        &=\int_s^t |\dot\rho_r|\,{\rm d}r
          =\int_s^t |\partial\F|(\rho_r)\,{\rm d}r \\
        &=\int_s^t \frac{|\partial\F|^2(\rho_r)}{|\partial\F|(\rho_r)}\,{\rm d}r \\
        &\le
          \frac{1}{\sqrt{2\kappa}}
          \int_s^t \frac{|\partial\F|^2(\rho_r)}{\sqrt{E(r)}}\,{\rm d}r \\
        &=\frac{1}{\sqrt{2\kappa}}
          \int_s^t \frac{-E'(r)}{\sqrt{E(r)}}\,{\rm d}r \\
        &=\sqrt{\frac{2}{\kappa}}\bigl(\sqrt{E(s)}-\sqrt{E(t)}\bigr).
\end{align*}
If $E$ vanishes on a subinterval, then $\rho_r$ is stationary there by \eqref{eq:full-ede}; the same estimate follows by applying the above argument on the part where $E>0$.
Since Wasserstein distance is bounded by metric length, \eqref{eq:w-between-times} follows.

Now let $t\to\infty$. From \eqref{eq:energy-rate-from-s}, $E(t)\to0$. From \eqref{eq:w-between-times}, the tail of the curve is Cauchy, so completeness gives a limit $\rho_\infty$. Lower semicontinuity yields
\[
        \F(\rho_\infty)
        \le \liminf_{t\to\infty}\F(\rho_t)
        =\F_\star,
\]
and hence $\rho_\infty\in\operatorname{Argmin}\F$. Sending the endpoint of \eqref{eq:w-between-times} to infinity with $s=t$ fixed gives the first inequality in \eqref{eq:w-to-limit-rate}; combining it with \eqref{eq:energy-rate-from-s} gives the second. Finally, \eqref{eq:w-to-argmin-rate} follows because $\rho_\infty$ belongs to $\operatorname{Argmin}\F$.
\end{proof}

\begin{corollary}[Quadratic-growth/error-bound consequence]
\label{cor:wpl-qg}
Assume that the hypotheses of Theorem~\ref{thm:wpl-flow-convergence} hold for the gradient flow starting from every $\rho\in\operatorname{Dom}\F$ with finite energy gap. Then the $\W$-PL inequality implies the static error bound
\begin{equation}
        \frac{\kappa}{2}
        \distW^2(\rho,\operatorname{Argmin}\F)
        \le
        \F(\rho)-\F_\star.
        \label{eq:qg-from-wpl}
\end{equation}
Equivalently,
\begin{equation}
        \distW(\rho,\operatorname{Argmin}\F)
        \le
        \sqrt{\frac{2}{\kappa}}
        \sqrt{\F(\rho)-\F_\star}.
        \label{eq:error-bound-from-wpl}
\end{equation}
\end{corollary}

\begin{proof}
Start the gradient flow at $\rho_0=\rho$ and apply \eqref{eq:w-to-argmin-rate} at time $t=0$.
\end{proof}

\begin{remark}[What the distance result says, and what it does not say]
The estimate \eqref{eq:w-to-argmin-rate} answers the natural question: under the standard gradient-flow identity, Wasserstein--PL does give a linear rate in Wasserstein distance to the minimizer selected by the flow. The exponent is $\kappa$, while the energy exponent is $2\kappa$, as expected after taking a square root.

This should not be confused with the stronger contraction estimate provided by positive Wasserstein geodesic convexity. If $\F$ is $\lambda$-geodesically convex and admits a minimizer $\rhostar$, the EVI theory typically gives
\[
        \W(\rho_t,\rhostar)
        \le e^{-\lambda t}\W(\rho_0,\rhostar),
\]
and even contraction between two solutions. Wasserstein--PL alone gives convergence to some minimizer and a distance-to-the-set estimate; it does not encode curvature along arbitrary geodesics and does not imply EVI contractivity.
\end{remark}

\begin{remark}[Nonunique minimizers]
The estimate concerns $\distW(\rho_t,\operatorname{Argmin}\F)$, or equivalently the distance to the limit minimizer $\rho_\infty$ selected by the flow. If the minimizer is not unique, Wasserstein--PL does not imply convergence to a prescribed minimizer. The finite-dimensional model $f(x,y)=x^2$ is a useful analogy: it satisfies PL, the gradient flow sends $(x_0,y_0)$ to $(0,y_0)$, and the distance to the solution set decays, but the distance to $(0,0)$ need not decay.
\end{remark}

\section{Positive Wasserstein geodesic convexity implies Wasserstein--PL}

\begin{definition}[Wasserstein strong convexity]
Let $\lambda>0$. The functional $\F$ is $\lambda$-geodesically convex, or $\lambda$-strongly convex in the Wasserstein sense, if for every constant-speed $\W$-geodesic $(\rho_s)_{s\in[0,1]}$,
\begin{equation}
        \F(\rho_s)
        \le (1-s)\F(\rho_0)+s\F(\rho_1)
        -\frac{\lambda}{2}s(1-s)\W^2(\rho_0,\rho_1).
        \label{eq:lambda-geodesic-convex}
\end{equation}
\end{definition}

\begin{proposition}[Strong convexity $\Rightarrow$ PL]\label{prop:strong-implies-pl}
Assume that $\F$ is $\lambda$-geodesically convex with $\lambda>0$, that $\F_\star$ is attained at some $\rhostar\in\Ptwo(M)$, and that the slope \eqref{eq:metric-slope} is finite at $\rho$. Then
\begin{equation}
        |\partial \F|^2(\rho)
        \ge 2\lambda\bigl(\F(\rho)-\F_\star\bigr).
        \label{eq:strong-implies-pl}
\end{equation}
Thus $\lambda$-Wasserstein strong convexity implies the Wasserstein--PL inequality with the same constant $\lambda$ in the convention \eqref{eq:wpl}.
\end{proposition}

\begin{proof}
Fix $\rho$ and a minimizer $\rhostar$. Put
\[
        g:=\F(\rho)-\F_\star,
        \qquad
        r:=\W(\rho,\rhostar),
        \qquad
        a:=|\partial\F|(\rho).
\]
If $r=0$ there is nothing to prove. Otherwise let $(\rho_s)_{s\in[0,1]}$ be a constant-speed geodesic from $\rho$ to $\rhostar$. By \eqref{eq:lambda-geodesic-convex},
\[
        \F(\rho_s)
        \le (1-s)\F(\rho)+s\F_\star
        -\frac{\lambda}{2}s(1-s)r^2.
\]
Hence
\[
        \F(\rho)-\F(\rho_s)
        \ge s g + \frac{\lambda}{2}s(1-s)r^2.
\]
Dividing by $\W(\rho,\rho_s)=sr$ and letting $s\downarrow0$ in the definition of the slope yields
\begin{equation}
        a\ge \frac{g}{r}+\frac{\lambda}{2}r.
        \label{eq:slope-geodesic-convex-bound}
\end{equation}
Equivalently,
\[
        g\le ar-\frac{\lambda}{2}r^2.
\]
Maximizing the right-hand side over $r\ge0$ gives
\[
        g\le \frac{a^2}{2\lambda},
\]
which is exactly \eqref{eq:strong-implies-pl}.
\end{proof}

\begin{remark}
This proof is the metric version of the elementary Euclidean argument: if $f$ is $\lambda$-strongly convex and $x_\star$ is a minimizer, then
\[
        f(x)-f(x_\star)
        \le \langle \nabla f(x),x-x_\star\rangle
          -\frac{\lambda}{2}|x-x_\star|^2
        \le \frac{1}{2\lambda}|\nabla f(x)|^2.
\]
The Wasserstein proof replaces $|x-x_\star|$ by $\W(\rho,\rhostar)$ and the Euclidean gradient norm by the metric slope.
\end{remark}

\begin{remark}[Strong convexity also gives static distance control]
If $\F$ is $\lambda$-geodesically convex with $\lambda>0$, then applying \eqref{eq:lambda-geodesic-convex} along a geodesic from a minimizer to $\rho$ also gives
\[
        \F(\rho)-\F_\star
        \ge \frac{\lambda}{2}
        \distW^2\bigl(\rho,\operatorname{Argmin}\F\bigr).
\]
Thus strong displacement convexity gives both PL and a static Wasserstein quadratic-growth bound. The PL theorem above shows that, for a genuine gradient flow, the distance rate can still be recovered even when geodesic convexity is absent.
\end{remark}

\section{The converse is false}

The PL inequality is first-order: it compares slope to height above the minimum. Strong convexity is second-order and asks for curvature along every geodesic. These are different requirements. The following example gives a genuine Wasserstein functional satisfying $\W$-PL while failing positive Wasserstein geodesic convexity.

\begin{example}[A nonconvex potential energy satisfying Wasserstein--PL]
Work on $\Ptwo(\R)$ and set
\begin{equation}
        V(x):=x^2+2\sin^2 x,
        \qquad
        \F(\rho):=\int_{\R} V(x)\dd\rho(x).
        \label{eq:counterexample-energy}
\end{equation}
Then $\F_\star=0$, attained uniquely at $\delta_0$.
The first variation is $\delta\F/\delta\rho=V$, so
\begin{equation}
        |\partial\F|^2(\rho)=\int_{\R}|V'(x)|^2\dd\rho(x)
        \label{eq:potential-slope}
\end{equation}
for every $\rho$ for which the right-hand side is finite. Now
\[
        V'(x)=2x+2\sin(2x),
        \qquad
        V''(x)=2+4\cos(2x).
\]
We first prove a pointwise PL inequality. Since $V$ is even and $V'$ is odd, it is enough to consider $r=|x|\ge0$.

If $0\le r\le \pi/2$, then $\sin(2r)\ge0$, hence
\[
        |V'(x)|\ge 2r,
        \qquad
        V(x)\le r^2+2r^2=3r^2,
\]
and therefore
\begin{equation}
        |V'(x)|^2\ge \frac{4}{3}V(x).
        \label{eq:pl-region-one}
\end{equation}
If $r\ge\pi/2$, then $2r-2>0$ and
\[
        |V'(x)|\ge 2r-2,
        \qquad
        V(x)\le r^2+2.
\]
The function $r\mapsto (2r-2)^2/(r^2+2)$ is increasing for $r>1$, so for $r\ge\pi/2$,
\begin{equation}
        |V'(x)|^2
        \ge
        \frac{(\pi-2)^2}{\pi^2/4+2}V(x).
        \label{eq:pl-region-two}
\end{equation}
Combining \eqref{eq:pl-region-one} and \eqref{eq:pl-region-two},
\begin{equation}
        |V'(x)|^2\ge c V(x)
        \quad\text{for all }x\in\R,
        \qquad
        c:=\min\left\{\frac43,\frac{(\pi-2)^2}{\pi^2/4+2}\right\}>0.
        \label{eq:pointwise-pl-counterexample}
\end{equation}
Integrating \eqref{eq:pointwise-pl-counterexample} against $\rho$ gives
\[
        |\partial\F|^2(\rho)
        =\int |V'|^2\dd\rho
        \ge c\int V\dd\rho
        = c\bigl(\F(\rho)-\F_\star\bigr).
\]
Thus $\F$ satisfies the Wasserstein--PL inequality \eqref{eq:wpl} with constant
\begin{equation}
        \kappa=\frac{c}{2}>0.
        \label{eq:counterexample-kappa}
\end{equation}

However, $\F$ is not $\lambda$-geodesically convex for any $\lambda>0$. Indeed, if it were, then restricting \eqref{eq:lambda-geodesic-convex} to Dirac masses would give, for all $x,y\in\R$ and $s\in[0,1]$,
\[
        V((1-s)x+sy)
        \le (1-s)V(x)+sV(y)-\frac{\lambda}{2}s(1-s)|x-y|^2,
\]
which is exactly $\lambda$-strong convexity of the scalar potential $V$. But
\[
        V''(\pi/2)=2+4\cos(\pi)=-2<0.
\]
So $V$ is not even convex, and positive Wasserstein geodesic convexity of $\F$ is impossible.
\end{example}

\begin{remark}
The example is not merely a finite-dimensional counterexample disguised as a Wasserstein one. The PL estimate holds after integration against every probability measure $\rho$, while failure of geodesic convexity is detected along the Wasserstein geodesics between Dirac masses. Thus the implication
\[
        \text{Wasserstein--PL} \quad\Longrightarrow\quad
        \text{positive Wasserstein geodesic convexity}
\]
fails on the full Wasserstein space.
\end{remark}

\section{A dictionary: functional inequalities as Wasserstein--PL inequalities}

This section explains how many entropy--entropy production inequalities are exactly Wasserstein--PL inequalities. The guiding principle is simple:
\[
        \boxed{
        \text{energy gap }\F(\rho)-\F_\star
        \quad\text{versus}\quad
        \text{Wasserstein slope squared }|\partial\F|^2(\rho)
        }
\]
A functional inequality of the form
\[
        2\kappa\bigl(\F(\rho)-\F_\star\bigr)
        \le |\partial\F|^2(\rho)
\]
is a Wasserstein--PL inequality by definition.

\subsection{Relative entropy and logarithmic Sobolev inequality}

Let $\pi$ be a probability measure on $\R^d$ or on a Riemannian manifold, with smooth positive density, and write $\rho=h\pi$. Consider the relative entropy
\begin{equation}
        \F(\rho)=\Ent_\pi(\rho)
        :=\int h\log h\dd\pi.
        \label{eq:relative-entropy}
\end{equation}
The first variation is $\log h$ up to an additive constant. Therefore
\begin{equation}
        |\partial\F|^2(\rho)
        =\int |\nabla\log h|^2 h\dd\pi
        =: I_\pi(\rho),
        \label{eq:fisher-information}
\end{equation}
the relative Fisher information. The $\W$-PL inequality for $\F$ is
\begin{equation}
        2\kappa\Ent_\pi(\rho)
        \le I_\pi(\rho),
        \label{eq:lsi-as-pl}
\end{equation}
which is exactly the logarithmic Sobolev inequality with constant $\kappa$.

The gradient flow is
\begin{equation}
        \partial_t\rho_t
        =\nabla\cdot\left(\rho_t\nabla\log\frac{\rho_t}{\pi}\right).
        \label{eq:fokker-planck-general}
\end{equation}
If $\pi=Z^{-1}e^{-V}\dd x$, this becomes the Fokker--Planck equation
\begin{equation}
        \partial_t\rho_t
        =\Delta\rho_t+\nabla\cdot(\rho_t\nabla V).
        \label{eq:fokker-planck}
\end{equation}
Under \eqref{eq:lsi-as-pl},
\[
        \Ent_\pi(\rho_t)
        \le e^{-2\kappa t}\Ent_\pi(\rho_0).
\]

\begin{example}[Gaussian case]
For the standard Gaussian measure $\gamma$ on $\R^d$,
\[
        2\Ent_\gamma(\rho)\le I_\gamma(\rho).
\]
Thus the entropy is a $\W$-PL functional with $\kappa=1$. The Ornstein--Uhlenbeck/Fokker--Planck flow has entropy decay $e^{-2t}$, and Theorem~\ref{thm:wpl-flow-convergence} also gives
\[
        \W(\rho_t,\gamma)
        \le \sqrt{2\Ent_\gamma(\rho_0)}e^{-t}.
\]
\end{example}


\subsection{Talagrand and quadratic-growth inequalities as distance consequences}

The distance estimate in Theorem~\ref{thm:wpl-flow-convergence} is not itself the Wasserstein--PL inequality; it is a consequence of Wasserstein--PL plus the energy dissipation identity. In the entropy case, it has the same shape as a Talagrand $T_2$ inequality. Indeed, if $\operatorname{Argmin}\Ent_\pi=\{\pi\}$ and the entropy gradient flow is global, then \eqref{eq:error-bound-from-wpl} gives
\begin{equation}
        \W^2(\rho,\pi)
        \le \frac{2}{\kappa}\Ent_\pi(\rho),
        \label{eq:t2-from-wpl-dynamic}
\end{equation}
whenever the flow can be started from $\rho$. For entropy, the classical Otto--Villani implication ``log-Sobolev implies Talagrand'' gives the same type of static estimate.

Thus there are three related but logically distinct inequalities:
\begin{align*}
        \text{Wasserstein--PL / LSI:}\quad
        &2\kappa\Ent_\pi(\rho)\le I_\pi(\rho),\\
        \text{Talagrand / transport inequality:}\quad
        &\W^2(\rho,\pi)\le \frac{2}{\kappa}\Ent_\pi(\rho),\\
        \text{flow distance rate:}\quad
        &\W(\rho_t,\pi)\le
          \sqrt{\frac{2}{\kappa}\Ent_\pi(\rho_0)}e^{-\kappa t}.
\end{align*}
The first inequality compares entropy with entropy production; the second compares transport distance with entropy; the third is a dynamical consequence along the gradient flow.

\subsection{\texorpdfstring{$\Phi$}{Phi}-entropy inequalities}

Let $\Phi:(0,\infty)\to\R$ be convex, with $\Phi(1)=\Phi'(1)=0$, and set
\begin{equation}
        \F_\Phi(\rho)=\int \Phi(h)\dd\pi,
        \qquad h=\frac{\dd\rho}{\dd\pi}.
        \label{eq:phi-entropy}
\end{equation}
The first variation is $\Phi'(h)$, hence
\begin{equation}
        |\partial \F_\Phi|^2(\rho)
        =\int h\,|\nabla\Phi'(h)|^2\dd\pi.
        \label{eq:phi-slope}
\end{equation}
The Wasserstein--PL inequality is therefore the $\Phi$-entropy production inequality
\begin{equation}
        2\kappa\int\Phi(h)\dd\pi
        \le
        \int h\,|\nabla\Phi'(h)|^2\dd\pi.
        \label{eq:phi-pl}
\end{equation}
This formula is a general template. The classical name of the inequality depends on the choice of $\Phi$.

\subsection{A weighted Poincare or chi-square production inequality}

Take
\[
        \Phi(r)=\frac12(r-1)^2.
\]
Then
\[
        \F(\rho)=\frac12\int (h-1)^2\dd\pi
        =\frac12\chi^2(\rho\mid\pi),
        \qquad
        |\partial\F|^2(\rho)=\int h|\nabla h|^2\dd\pi.
\]
The corresponding Wasserstein--PL inequality is
\begin{equation}
        \kappa\int (h-1)^2\dd\pi
        \le \int h|\nabla h|^2\dd\pi.
        \label{eq:weighted-poincare-pl}
\end{equation}
This is a weighted Poincare-type inequality, or a chi-square entropy production inequality, with the Wasserstein mobility $h$ on the right-hand side. It should not be confused with the standard Poincare inequality
\[
        \kappa\int (h-1)^2\dd\pi\le \int |\nabla h|^2\dd\pi,
\]
although the two are directly comparable on classes where $h$ is bounded below or above.

The associated Wasserstein gradient flow is nonlinear:
\begin{equation}
        \partial_t\rho_t=\nabla\cdot(\rho_t\nabla h_t),
        \qquad h_t=\frac{\dd\rho_t}{\dd\pi}.
        \label{eq:chi-square-flow}
\end{equation}

\subsection{Power entropies, Beckner-type inequalities, and porous-medium structure}

For $m>1$, define the normalized power entropy
\begin{equation}
        \Phi_m(r)=\frac{r^m-mr+m-1}{m-1}.
        \label{eq:power-phi}
\end{equation}
Then $\Phi_m(1)=\Phi_m'(1)=0$ and
\[
        \Phi_m'(r)=\frac{m}{m-1}(r^{m-1}-1).
\]
Using \eqref{eq:phi-slope},
\begin{align}
        |\partial\F_{\Phi_m}|^2(\rho)
        &=m^2\int h^{2m-3}|\nabla h|^2\dd\pi \notag\\
        &=\frac{4m^2}{(2m-1)^2}
          \int \left|\nabla h^{m-1/2}\right|^2\dd\pi.
        \label{eq:power-slope}
\end{align}
Thus the $\W$-PL inequality is
\begin{equation}
        2\kappa\int \frac{h^m-mh+m-1}{m-1}\dd\pi
        \le
        \frac{4m^2}{(2m-1)^2}
        \int \left|\nabla h^{m-1/2}\right|^2\dd\pi.
        \label{eq:power-pl}
\end{equation}
Depending on the reference measure and the range of exponents, inequalities of this type appear as Beckner, convex Sobolev, generalized Sobolev, or interpolation inequalities. In the limit $m\downarrow1$, \eqref{eq:power-pl} formally recovers the logarithmic Sobolev inequality.

The free diffusion energy
\[
        \rho\mapsto \frac{1}{m-1}\int_{\R^d}\rho^m\dd x
\]
has Wasserstein gradient flow
\[
        \partial_t\rho_t=\Delta(\rho_t^m),
\]
the porous-medium equation. In confined or self-similar variables, the corresponding relative entropy--production inequality gives exponential convergence to the Barenblatt equilibrium and is closely related to sharp Gagliardo--Nirenberg--Sobolev inequalities.

\subsection{Nonlinear Fokker--Planck and granular-media free energies}

Consider a free energy on $\Ptwo(\R^d)$ of the form
\begin{equation}
        \F(\rho)
        =\int \rho\log\rho\dd x
         +\int V(x)\dd\rho(x)
         +\frac12\iint W(x-y)\dd\rho(x)\dd\rho(y).
        \label{eq:granular-free-energy}
\end{equation}
The first variation is
\[
        \frac{\delta\F}{\delta\rho}
        =\log\rho+1+V+W*\rho,
\]
so the Wasserstein slope squared is
\begin{equation}
        |\partial\F|^2(\rho)
        =\int \left|\nabla\log\rho+\nabla V+\nabla W*\rho\right|^2\dd\rho.
        \label{eq:granular-dissipation}
\end{equation}
The $\W$-PL inequality is the entropy--entropy production estimate
\begin{equation}
        2\kappa\bigl(\F(\rho)-\F_\star\bigr)
        \le
        \int \left|\nabla\log\rho+\nabla V+\nabla W*\rho\right|^2\dd\rho.
        \label{eq:granular-pl}
\end{equation}
The associated gradient flow is the aggregation--diffusion equation
\begin{equation}
        \partial_t\rho_t
        =\nabla\cdot\left(\rho_t\nabla(\log\rho_t+V+W*\rho_t)\right).
        \label{eq:aggregation-diffusion}
\end{equation}
When the energy is positively displacement convex, \eqref{eq:granular-pl} follows from Proposition~\ref{prop:strong-implies-pl}. In some nonconvex settings, an inequality of the form \eqref{eq:granular-pl} may still hold and then gives exponential decay of the free energy without geodesic convexity.

\subsection{HWI as a route to PL for entropy}

For entropy relative to a measure $\pi$ satisfying a positive curvature condition, the HWI inequality has the schematic form
\begin{equation}
        \Ent_\pi(\rho)
        \le
        \W(\rho,\pi)\sqrt{I_\pi(\rho)}
        -\frac{K}{2}\W^2(\rho,\pi),
        \qquad K>0.
        \label{eq:hwi}
\end{equation}
Optimizing the right-hand side over the distance variable gives
\begin{equation}
        \Ent_\pi(\rho)\le \frac{1}{2K}I_\pi(\rho),
        \label{eq:hwi-to-lsi}
\end{equation}
that is, the logarithmic Sobolev inequality and hence the Wasserstein--PL inequality for entropy. Once this PL inequality is available, Theorem~\ref{thm:wpl-flow-convergence} gives the corresponding flow estimate
\[
        \W(\rho_t,\pi)
        \le \sqrt{\frac{2}{K}\Ent_\pi(\rho_0)}e^{-Kt}
\]
for the entropy gradient flow. Thus HWI is not itself the PL inequality, but it is a standard mechanism producing PL; the PL inequality plus the gradient-flow identity then gives the distance rate.

\subsection{Summary table}

\begin{center}
\renewcommand{\arraystretch}{1.35}
\footnotesize
\begin{tabular}{>{\raggedright\arraybackslash}p{0.22\textwidth} >{\raggedright\arraybackslash}p{0.25\textwidth} >{\raggedright\arraybackslash}p{0.39\textwidth}}
\toprule
Energy $\F(\rho)$ & Slope squared $|\partial\F|^2$ & $\W$-PL inequality / usual name \\
\midrule
$\Ent_\pi(\rho)=\int h\log h\dd\pi$ & $I_\pi(\rho)=\int |\nabla\log h|^2h\dd\pi$ & $2\kappa\Ent_\pi\le I_\pi$; logarithmic Sobolev inequality. \\
$\frac12\chi^2(\rho\mid\pi)$ & $\int h|\nabla h|^2\dd\pi$ & Weighted Poincare / chi-square entropy production inequality. \\
$\int \Phi(h)\dd\pi$ & $\int h|\nabla\Phi'(h)|^2\dd\pi$ & $\Phi$-entropy production inequality. \\
Power entropy $\int \Phi_m(h)\dd\pi$ & $\frac{4m^2}{(2m-1)^2}\int |\nabla h^{m-1/2}|^2\dd\pi$ & Beckner, convex Sobolev, generalized Sobolev, or interpolation inequality. \\
Confined porous-medium entropy & $\int \rho|\nabla\psi_\rho|^2\dd x$, $\psi_\rho=\frac{m}{m-1}\rho^{m-1}+V$ & Entropy--entropy production inequality; often equivalent to Gagliardo--Nirenberg--Sobolev inequalities in self-similar variables. \\
Granular-media free energy \eqref{eq:granular-free-energy} & Dissipation \eqref{eq:granular-dissipation} & Entropy--entropy production inequality for aggregation--diffusion. \\
\bottomrule
\end{tabular}
\normalsize
\end{center}


\section{Takeaways}

\begin{enumerate}[label=(\roman*)]
\item A Wasserstein--PL inequality is the metric-slope inequality
\[
        |\partial\F|^2\ge 2\kappa(\F-\F_\star).
\]
Together with the Wasserstein energy dissipation identity, it gives the energy rate
\[
        \F(\rho_t)-\F_\star\le e^{-2\kappa t}(\F(\rho_0)-\F_\star).
\]

\item The same assumptions also give a distance rate. The key estimate is the finite-length bound
\[
        \operatorname{Length}_{\W}\bigl((\rho_s)_{s\ge t}\bigr)
        \le \sqrt{\frac{2}{\kappa}}
        \sqrt{\F(\rho_t)-\F_\star}.
\]
Hence, under lower semicontinuity and completeness assumptions,
\[
        \distW(\rho_t,\operatorname{Argmin}\F)
        \le
        \sqrt{\frac{2}{\kappa}}
        \sqrt{\F(\rho_0)-\F_\star}\,e^{-\kappa t}.
\]
The exponent is half the energy exponent because distance is controlled by the square root of the energy gap.

\item Positive Wasserstein geodesic convexity is sufficient for Wasserstein--PL and gives stronger EVI/contractivity information. Wasserstein--PL by itself gives convergence to the minimizer selected by the flow, but not curvature along arbitrary Wasserstein geodesics.

\item The converse is false. The potential energy $\F(\rho)=\int (x^2+2\sin^2x)\dd\rho(x)$ satisfies Wasserstein--PL because the pointwise scalar PL inequality $|V'|^2\ge cV$ holds, but it is not positively Wasserstein geodesically convex because $V''(\pi/2)<0$.

\item Many standard functional inequalities are simply Wasserstein--PL inequalities written for specific choices of $\F$. The most important example is the logarithmic Sobolev inequality, which is the Wasserstein--PL inequality for relative entropy. Talagrand-type inequalities are not PL inequalities, but they are closely related distance/error-bound consequences.
\end{enumerate}

\begin{thebibliography}{99}

\bibitem{AGS}
L. Ambrosio, N. Gigli, and G. Savar\'e.
\newblock \emph{Gradient Flows in Metric Spaces and in the Space of Probability Measures}.
\newblock Birkh\"auser, second edition, 2008.

\bibitem{JKO}
R. Jordan, D. Kinderlehrer, and F. Otto.
\newblock The variational formulation of the Fokker--Planck equation.
\newblock \emph{SIAM Journal on Mathematical Analysis}, 29(1):1--17, 1998.

\bibitem{OttoVillani}
F. Otto and C. Villani.
\newblock Generalization of an inequality by Talagrand, and links with the logarithmic Sobolev inequality.
\newblock \emph{Journal of Functional Analysis}, 173(2):361--400, 2000.

\bibitem{GigliLedoux}
N. Gigli and M. Ledoux.
\newblock From log Sobolev to Talagrand: a quick proof.
\newblock \emph{Discrete and Continuous Dynamical Systems}, 33(5):1927--1935, 2013.

\bibitem{CMV}
J. A. Carrillo, R. J. McCann, and C. Villani.
\newblock Kinetic equilibration rates for granular media and related equations: entropy dissipation and mass transportation estimates.
\newblock \emph{Revista Matem\'atica Iberoamericana}, 19(3):971--1018, 2003.

\bibitem{CarrilloToscani}
J. A. Carrillo and G. Toscani.
\newblock Asymptotic $L^1$-decay of solutions of the porous medium equation to self-similarity.
\newblock \emph{Indiana University Mathematics Journal}, 49(1):113--142, 2000.


\bibitem{GuptaBalakrishnanRamdas}
C. Gupta, S. Balakrishnan, and A. Ramdas.
\newblock Path length bounds for gradient descent and flow.
\newblock \emph{Journal of Machine Learning Research}, 22(68):1--62, 2021.

\bibitem{BolteDaniilidisLeyMazet}
J. Bolte, A. Daniilidis, A. Ley, and L. Mazet.
\newblock Characterizations of {\L}ojasiewicz inequalities: subgradient flows, talweg, convexity.
\newblock \emph{Transactions of the American Mathematical Society}, 362(6):3319--3363, 2010.

\bibitem{LiaoLiZhang}
F. Y. Liao, L. Li, and Y. Zhang.
\newblock Error bounds, PL condition, and quadratic growth for weakly convex functions, and linear convergences of proximal point methods.
\newblock \emph{Proceedings of Machine Learning Research}, 242:2029--2057, 2024.

\bibitem{KarimiNutiniSchmidt}
H. Karimi, J. Nutini, and M. Schmidt.
\newblock Linear convergence of gradient and proximal-gradient methods under the Polyak--{\L}ojasiewicz condition.
\newblock In \emph{ECML PKDD}, pages 795--811, 2016.

\bibitem{BoufadeneVialard}
S. Boufad\`ene and F.-X. Vialard.
\newblock On the global convergence of Wasserstein gradient flow of the Coulomb discrepancy.
\newblock \emph{SIAM Journal on Mathematical Analysis}, 2025.

\end{thebibliography}

\end{document}
